% Chapter 3, section 11: One worker at a time
% Write in plain LaTeX. Code of a guide step: \guide{<step>}; commands: \begin{shell}...\end{shell};
% a screenshot: \screenshot{file}{caption}; a gap in the draft: \placeholder{...}.
% See book/latex/queen.sty for the full list.
%
% Draft by Claude, 2026-10-24. The code is cut from examples/tutorials/js/book-use-cases.mjs, which
% checks every result quoted here against a 2.6.0 node.

\section{One worker at a time}

The nightly report must be built once, but every replica of the service runs the same scheduler.
This is leader election, or a singleton job.

A lock with a lifetime: the worker that gets it builds the report, and the others skip. The report
is committed with \texttt{.guard(lock)}: the transaction commits only while the lock is still this
worker's. The day in the \texttt{transactionId} is a second fence.

\queenfig{singleton}{One worker builds the report; the guard keeps a late one out.}

\guide{uc-singleton}

Three workers started at once: one built the report and two skipped. There was one report.

\why{a lock expires without telling its holder. A worker paused past the lifetime keeps running,
while another one takes the lock. The guard keeps the late worker's work out: its commit rolls
back.}

One thing to know: the lock renews itself in the background while the worker holds it. For long
work, pass \texttt{lock.signal} to what the work waits on, so it stops when the lock is lost. For
``at most three at a time'', \texttt{queen.semaphore(name, 3)} works the same way.
